\documentclass[border=10pt]{standalone}
\usepackage{ctex}
\usepackage[european]{circuitikz}
\begin{document}
\begin{circuitikz}
% 1. 下降沿 JK 触发器：c2=1 (clock), n2=1 (negation bubble), dot on notQ (n4=1)
\draw (0,0) node[flipflop, flipflop def={t1=J, t2={\footnotesize C1}, t3=K, t6=Q, t4=Q, c2=1, n2=1, n4=1}] (ff1) {};
\node[above] at (ff1.north) {JK (falling, Qbar with bubble)};

% 2. 上升沿 D 触发器：c2=1 (clock), 无 bubble
\draw (4,0) node[flipflop, flipflop def={t1=1D, t2={\footnotesize C1}, t6=Q, t4=Q, c2=1, n4=1}] (ff2) {};
\node[above] at (ff2.north) {D (rising, Qbar with bubble)};

% 3. 74LS290 / 74LS161 用 dipchip 或者 custom multiplexer?
% circuitikz 提供了 dipchip
% 还有 muxdemux / chips
% 测试 dipchip:
\draw (8,0) node[dipchip, num pins=12, hide numbers, no marks] (c1) {};

% 测试 spdt 的不同方向和切换
% spdt 的触点：
\draw (0,-4) node[spdt] (sw1) {};
\draw (3,-4) node[spdt, xscale=-1] (sw2) {};

\end{circuitikz}
\end{document}
